
\newpage
\thispagestyle{empty} 
\setcounter{page}{2}
\addcontentsline{toc}{section}{Abstract} 
\section*{\centering Abstract}
\noindent 
Automation is progressively displacing skilled technical labor, however, current robotic systems remain limited in replicating complex tasks, leaving a persistent need for human operators. At the same time, such skills are difficult to acquire, requiring prolonged training and the development of muscle memory. This creates a gap between the continued demand for skilled operators and the accessibility of training pathways.

\noindent
This paper presents a design-based research study investigating Extended Reality (XR) as a medium for lowering the learning barrier and increasing engagement in complex skill acquisition, using First Person View (FPV) quadcopter piloting as a representative case domain.

\noindent
This study is structured around three primary objectives: (1) the simulation of FPV quadcopter flight physics. (2) the design of a human-AI co-training framework (3) the interaction instruments for a high-fidelity XR prototyping. In particular PID and RL (Reinforcement Learning) are used in the co-training AI, and Geographic Information System (GIS) data integrated to support the environmental representation.

\noindent
The contribution of this work is a structured design framework that links physical simulation, AI-assisted co-training, and XR-based interaction design, aimed at reducing the learning curve and improving accessibility to complex skill training.

\vspace{0.5in} 
\textit{Keywords:} 
Extended Reality, 
Digital Twin, 
FPV Drone, 
Human-AI Co-training,
PID Control,
Reinforcement Learning
